
\subsection{Código de la Pregunta a): función \texttt{calcular\_TVI}}
\label{anexo:codigo_a}

\definecolor{gray}{gray}{0.5}

\begin{lstlisting}
import numpy as np

BETA = 0.05
VST = 0.9
T_FALLA = 0.0
T_DESPEJE = 0.1
T_END = 20.0
TOL = 1e-6

_trapz = getattr(np, "trapezoid", None) or np.trapz


def envolvente(t, beta=BETA, vst=VST, t_end=T_END):
    """Envolventes inferior y superior segun Wildenhues et al."""
    t = np.asarray(t, dtype=float)
    t_low = ((t / t_end) * np.exp(t / t_end)) ** beta / np.exp(beta) * vst
    t_upp = 2.0 - t_low
    return t_low, t_upp


def calcular_TVI(app, nombre_linea):
    """Retorna [[barra, TVI], ...] con TVI != 0, [] si no hay
    violaciones o False si la simulacion no converge."""
    try:
        # 1. Preparar eventos (se reutilizan los del caso de estudio)
        evt = app.GetFromStudyCase("IntEvt")
        linea = app.GetCalcRelevantObjects(nombre_linea + ".ElmLne")[0]
        falla = evt.GetContents("ShcLineEvent*")[0]
        disparo = evt.GetContents("TripLineEvent*")[0]
        for e in evt.GetContents():
            e.outserv = 1
        falla.outserv = 0
        disparo.outserv = 0

        # 2. Cortocircuito trifasico franco al 50 % en t = 0
        falla.p_target = linea
        falla.time = T_FALLA
        falla.i_shc = 0
        falla.R_f = 0.0
        falla.X_f = 0.0
        falla.fshcloc = 50.0

        # 3. Apertura de la linea en t = 0.1 s (evento, no outserv)
        disparo.p_target = linea
        disparo.time = T_DESPEJE
        disparo.i_switch = 0

        # 4. Condiciones iniciales y simulacion RMS hasta 20 s
        inc = app.GetFromStudyCase("ComInc")
        sim = app.GetFromStudyCase("ComSim")
        inc.iopt_sim = "rms"
        if inc.Execute() != 0:
            return False
        sim.tstop = T_END
        if sim.Execute() != 0:
            return False

        # 5. Extraer tensiones en p.u. de las barras
        res = inc.p_resvar
        res.Load()
        n = res.GetNumberOfRows()
        t = np.array([res.GetValue(i, -1)[1] for i in range(n)])
        barras = app.GetCalcRelevantObjects("*.ElmTerm")
        V = {}
        for b in barras:
            col = res.FindColumn(b, "m:u")
            if col < 0:
                continue
            V[b.loc_name] = np.array(
                [res.GetValue(i, col)[1] for i in range(n)])
        res.Release()
    except Exception:
        return False

    # 6. Envolventes y TVI por barra (desde el despeje hasta t_end)
    mask = t >= T_DESPEJE
    tc = t[mask]
    t_low, t_upp = envolvente(tc)

    salida = []
    for nombre, v in V.items():
        vc = v[mask]
        viol = np.maximum(t_low - vc, 0.0) + np.maximum(vc - t_upp, 0.0)
        tvi = float(_trapz(viol, tc))
        if tvi > TOL:
            salida.append([nombre, tvi])
    return salida
\end{lstlisting}